% section: 発展的な解法

\section{発展的な解法}\label{節：発展的な解法}
倍率が変わる場合,整数に丸める場合,更新に使う項が増える場合へ進もう.
追う量を変え,更新を求めて元へ戻す流れは共通している.\footnote{ここで扱う\bookterm{recurrence}{漸化式}は,一部の難関大でも出題される.
  基本的に誘導がつくので,何を簡単にしているかを読むことが大切である.}
最初は変わる倍率を除き,\bookterm{ratio}{階比}型と\bookterm{difference}{階差}型を組み合わせる.
数表の横と縦で階差を取る広がりは,付録の\sectionref*{節：二方向の階差}で扱う.
\subsection{変数係数の一次型}\label{節：変数係数の一次型}
倍率と付加項が添字に依存する式を,\sectionref*{節：階比型}の積と\sectionref*{節：階差等比複合型}の引き算で解く.

\subsubsection{例題：倍率をそろえる}
まず
\BookMathLayout{\[
    a_1=2,\qquad
    a_{n+1}=\frac{n+2}{n}a_n+(n+1)(n+2)
    \quad(n\ge1)
  \]}{\[
    \begin{aligned}
      a_1&=2,\\
      a_{n+1}&=\frac{n+2}{n}a_n+(n+1)(n+2)\\
      &\quad(n\ge1)
    \end{aligned}
  \]}
を考える.
倍率の分母$n$と分子$n+2$の間には$n+1$がある.
$n(n+1)$と一つ先の$(n+1)(n+2)$を比べれば,共通する$n+1$が消えて元の倍率を作れる.
\[
  \frac{(n+1)(n+2)}{n(n+1)}=\frac{n+2}{n}
\]
から,\textbf{隣り合う比が元の倍率になる量}$h_n=n(n+1)$を選ぶ(\sectionref*{節：補助数列・正規化}).
$n\ge1$で非零なので,$b_n=\dfrac{a_n}{h_n}$と置いて割ると,
\[
  \frac{a_{n+1}}{(n+1)(n+2)}
  =\frac{a_n}{n(n+1)}+1
\]
となる.従って$b_1=1$, $b_{n+1}=b_n+1$より$b_n=n$であり,
\[
  a_n=n^2(n+1)
\]
を得る.

\bookstep{初項を変える}
初項を一般の$a_1=s$に変えると,$b_1=\dfrac{s}{2}$となる.
同じ更新から$b_n=\dfrac{s}{2}+n-1$なので,
\[
  a_n=n(n+1)\left(n+\frac{s-2}{2}\right)
\]
となる.

\subsubsection{一般の式への拡張}
$a_{n+1}=p_na_n+q_n$ $(n\ge1)$に対して,
$h_1=1$, $h_{n+1}=p_nh_n$とする.
すべての$p_n$が非零なら,$h_n$も非零であり,
\[
  h_n=\prod_{k=1}^{n-1}p_k
\]
である.\footnote{空積は$1$とする.}元の式を$h_{n+1}$で割ると,
\[
  b_n=\frac{a_n}{h_n},\qquad
  b_{n+1}-b_n=\frac{q_n}{h_{n+1}}
\]
となる.$b_1=a_1$から\bookterm{difference}{階差}を足すと,
\[
  b_n=a_1+\sum_{k=1}^{n-1}\frac{q_k}{h_{k+1}}
\]
となり,$a_n=h_nb_n$へ戻して,
\[
  a_n=h_n\left(a_1+\sum_{k=1}^{n-1}\frac{q_k}{h_{k+1}}\right)
\]
を得る.\footnote{先ほどは定数倍して$h_n=n(n+1)$を選んだが,
  その場合の$b_1$を正しく合わせれば同じ答えになる.}
\begin{bookpoint}
  \textbf{\bookterm{normalize}{正規化}の役割}\par
  倍率を取り除くと,残った付加項を足す問題になる.
  積・和の表示を得ることと,最後まで簡単な式で計算できることは別である.
\end{bookpoint}
もし$p_j=0$なら,$a_{j+1}=q_j$はそれ以前の項によらず決まる.
零となる更新で区間を分け,非零の倍率が続く範囲で割る方法を使う.

\subsubsection{別解：同じ更新を引く}
同じ例で,正規化から見つけた$v_n=n^2(n+1)$は元の更新を満たす.
今度は$d_n=a_n-v_n$と引けば,
\[
  d_{n+1}=\frac{n+2}{n}d_n
\]
となり,\bookterm{ratio}{階比}型へ戻る.初項$a_1=s$なら$d_1=s-2$であり,
\[
  d_n=\frac{s-2}{2}n(n+1),\qquad a_n=v_n+d_n
\]
から同じ答えを得る.一つの解$v_n$と\bookterm{homogeneous}{同次}の解の定数倍に分かれた(\sectionref*{節：階差等比複合型}).

\begin{bookpoint}
  \textbf{引くか,割るか}\par
  隣り合う比から割る因子を探すか,同じ更新をする量を引く.
  どちらも,付加項と倍率を分ける操作である.
\end{bookpoint}
操作だけの比較は\sectionref*{節：操作を選ぶ練習},
\bookterm{general}{一般項}まで求める練習は\bookproblemref{practice-connections}{1}{接続演習問1}へ進もう.

\subsection{ガウス記号を含む式}\label{節：ガウス記号を含む式}
\bookterm{floor}{ガウス記号}を含む\bookterm{recurrence}{漸化式}も,既知の更新へ帰着させて解く.
手掛かりは整数性,不等式,余り,偶奇,数え上げである.
\textbf{記号が式の中で何をしているか}から十の使い方を整理し,それぞれに例題を置く.
群と数え上げのように,一つの問題に複数の見方が重なることもある.

\subsubsection{基本性質と使い方}\label{節：基本性質と使い方}
\bookterm{floor}{ガウス記号}$\lfloor x\rfloor$または$[x]$は,実数$x$を超えない最大の整数を表し,\bookterm{floor-function}{床関数}と呼ばれる.\footnote{対になる記号として,天井関数$\lceil x\rceil$は$x$以上の最小の整数を表す.
  整数$k$について$\lceil x\rceil=k\iff k-1<x\le k$である.以下の例題では床関数を使う.}
整数$k$について,
\[
  \lfloor x\rfloor=k\iff k\le x<k+1
\]
である.\,$x\ge0$では小数点以下の切り捨てに一致するが,負数では$\lfloor-1.5\rfloor=-2$となる.

\BookFigFloorGraph

\bookterm{floor-function}{床関数}が\textbf{項の値,添字,付加項のどれに作用するか}を確かめ,次の一覧を手掛かりに読む.
\begin{enumerate}[label=(\arabic*)]
  \item 整数性・不等式で整数部分を確定する.
  \item 同じ値を取る区間を群としてまとめる.
  \item 余りを取り出し,有限個の値と周期を調べる.
  \item \bookterm{floor-function}{床関数}の差で境界・条件・繰り上がりを検出する.
  \item 丸め誤差と係数の関係を使う.
  \item 整数への丸めで\bookterm{fixed}{不動点}と停止を調べる.
  \item 添字を二つに分割する.
  \item 商と余りで整数の桁を読む.
  \item 小数部分を残し,回転として読む.
  \item 条件を満たす整数や格子点の個数を数える.
\end{enumerate}
桁を読む例題では,添字$n$の商と余りを使う.
\bookterm{general}{一般項}だけでなく,周期,止まる値,数えている対象も調べる.

整数$k$について$\lfloor x+k\rfloor=\lfloor x\rfloor+k$だが,実数同士では和に分配できない.
数項から得た予想は,不等式や帰納法で確かめる.

\subsubsection{整数部分を確定する}\label{節：整数部分を確定する}
\bookterm{floor}{ガウス記号}の中の整数部分を探す.
整数$k$と実数$x$に対して,
\[
  \lfloor k+x\rfloor=k+\lfloor x\rfloor
\]
が成り立つ.整数$a_n$なら$\lfloor a_n+\sqrt3\rfloor=a_n+1$である.
先に,初項から整数性が保たれることを確かめる.

中身を$k\le X<k+1$に挟んでもよい.平方根なら隣り合う平方数を使う.
\textbf{\bookterm{floor-function}{床関数}の値を確定し,既知の\bookterm{recurrence}{漸化式}へ戻す.}

\noindent \bookblocklink{example-gauss-role-1}{problem}{answer}{【例題】}\par\nobreak
\begin{enumerate}[label=(\arabic*),parsep=3pt,beginpenalty=10000,format=\bookitemlink{example-gauss-role-1}{problem}{answer}]
  \item \bookterm{sequence}{数列}を次のように定める.\bookterm{general}{一般項}を求めよ.
    \BookMathLayout{            \[
        a_1=1,\qquad
        a_{n+1}=\left\lfloor\sqrt{a_n^2+4a_n+5}\right\rfloor
        \quad(n\ge1).
      \]}{            \[
        \begin{aligned}
          a_1&=1,\\
          a_{n+1}&=\left\lfloor\sqrt{a_n^2+4a_n+5}\right\rfloor
          \\              &\quad(n\ge1).
        \end{aligned}
      \]}
\end{enumerate}

\noindent \bookblocklink{example-gauss-role-1}{hint}{problem}{【方針】}\par\nobreak
\begin{enumerate}[label=(\arabic*),parsep=3pt,beginpenalty=10000,format=\bookitemlink{example-gauss-role-1}{hint}{problem}]
  \item 整数$a_n$に対して,平方根の中身を$(a_n+2)^2$と$(a_n+3)^2$で挟みます.
    まず各項が非負整数であることを確かめましょう.
\end{enumerate}

\noindent \bookblocklink{example-gauss-role-1}{answer}{problem}{【解答】}\par\nobreak
\begin{enumerate}[label=(\arabic*),parsep=3pt,beginpenalty=10000,format=\bookitemlink{example-gauss-role-1}{answer}{problem}]
  \item \[
      a_n=2n-1\quad(n\ge1).
    \]
\end{enumerate}

\noindent \bookblocklink{example-gauss-role-1}{solution}{problem}{【解法】}\par\nobreak
\begin{enumerate}[label=(\arabic*),parsep=3pt,beginpenalty=10000,format=\bookitemlink{example-gauss-role-1}{solution}{problem}]
  \item 初項と更新から,帰納的に各項は非負整数となる.
    非負整数$x$について,
    \[
      (x+2)^2\le x^2+4x+5<(x+3)^2
    \]
    である(左の差は1,右の差は$2x+4>0$).平方根を取ると,
    \[
      x+2\le\sqrt{x^2+4x+5}<x+3
    \]
    より$a_{n+1}=a_n+2$.\sectionref{節：等差型}に帰着し,$a_n=2n-1$を得る.
\end{enumerate}

\subsubsection{同じ値を群にまとめる}\label{節：同じ値を群にまとめる}
\bookterm{floor-function}{床関数}が同じ値を取る項を一つの群にまとめる.
まず,一般項に\bookterm{floor}{ガウス記号}がある\bookterm{sequence}{数列}を見よう.

\noindent\phantomsection\label{証明：ガウス型の群数列}\textbf{【命題1】}\par\nobreak
\begin{smallquote}
  一般項が
  \[
    a_n=\lfloor \sqrt{n} \rfloor\quad(n\ge1)
  \]
  と定義される\bookterm{sequence}{数列}$a_n$は群\bookterm{sequence}{数列}を成す.
\end{smallquote}
\noindent \textbf{【証明1】}\par\nobreak
\begin{smallquote}
  正整数$k$について,$a_n=k$となる番号$n$は,
  \[
    k\le\sqrt{n}<k+1
  \]
  $n,k\ge1$から,
  \begin{align*}
    & k^2\le n <(k+1)^2                      \\
    & \therefore\quad k^2\le n \le (k+1)^2-1
  \end{align*}
  この範囲にある整数$n$の個数は
  \[
    \{(k+1)^2-1\}-k^2+1=2k+1
  \]
  である.値$k$が$n=k^2$から$2k+1$個並び,次の群では$k+1$となる. \qed
\end{smallquote}

実際の並びも,
\[
  \underbrace{1,1,1}_{3\text{個}},\,
  \underbrace{2,2,2,2,2}_{5\text{個}},\,
  \underbrace{3,3,3,3,3,3,3}_{7\text{個}},\,4,\ldots
\]
となる.

\bookterm{difference}{階差}なら\sectionref{節：階差型}として足す.
\textbf{群の境界と,和の上端が最後の群のどこにあるか}を区別する.

\noindent \bookblocklink{example-gauss-role-2}{problem}{answer}{【例題】}\par\nobreak
\begin{enumerate}[label=(\arabic*),parsep=3pt,beginpenalty=10000,format=\bookitemlink{example-gauss-role-2}{problem}{answer}]
  \item \bookterm{sequence}{数列}を次のように定める.\bookterm{general}{一般項}を求めよ.
    \[
      a_1=2,\qquad
      a_{n+1}=a_n+\lfloor\sqrt{n+2}\rfloor
      \quad(n\ge1).
    \]
    必要ならば,\,$m=\lfloor\sqrt{n+1}\rfloor$を用いてよい.
\end{enumerate}

\noindent \bookblocklink{example-gauss-role-2}{hint}{problem}{【方針】}\par\nobreak
\begin{enumerate}[label=(\arabic*),parsep=3pt,beginpenalty=10000,format=\bookitemlink{example-gauss-role-2}{hint}{problem}]
  \item \bookterm{difference}{階差}を足し,\,$\lfloor\sqrt k\rfloor$の和にします.
    \,$k=1,2$で値がともに1であることを使い,和の範囲をそろえましょう.
\end{enumerate}

\noindent \bookblocklink{example-gauss-role-2}{answer}{problem}{【解答】}\par\nobreak
\begin{enumerate}[label=(\arabic*),parsep=3pt,beginpenalty=10000,format=\bookitemlink{example-gauss-role-2}{answer}{problem}]
  \item \,$m=\lfloor\sqrt{n+1}\rfloor$とおくと,
    \[
      a_n=m(n+2)-\frac{m(m+1)(2m+1)}6.
    \]
\end{enumerate}

\noindent \bookblocklink{example-gauss-role-2}{solution}{problem}{【解法】}\par\nobreak
\begin{enumerate}[label=(\arabic*),parsep=3pt,beginpenalty=10000,format=\bookitemlink{example-gauss-role-2}{solution}{problem}]
  \item \bookterm{difference}{階差}を足すと,
    \[
      a_n=2+\sum_{k=3}^{n+1}\lfloor\sqrt k\rfloor
      =\sum_{k=1}^{n+1}\lfloor\sqrt k\rfloor.
    \]\footnote{\,$n=1$では最初の和を空の和として0とする.}
    \,$N=n+1,\ m=\lfloor\sqrt N\rfloor$とおく.
    項番号$k$で足し,第$m-1$群までは全部,第$m$群は$N$までを使う.
    \[
      \begin{array}{c|c|c}
        \text{群}&\text{足す項番号}&\text{個数}\\ \hline
        j<m&\substack{j^2\le k\\ k<(j+1)^2}&2j+1\\
        m&\substack{m^2\le k\\ k\le N}&N-m^2+1
      \end{array}
    \]
    最後の群は両端を含むので$N-m^2+1$個.値と個数を掛けて足すと,
    \begin{align*}
      a_n&=\sum_{j=1}^{m-1}j(2j+1)+m(N-m^2+1)\\
      &=\frac{(m-1)m(4m+1)}6\\
      &\quad+m(N-m^2+1)\\
      &=m(N+1)-\frac{m(m+1)(2m+1)}6.
    \end{align*}
    \,$N=n+1$を戻す.\,$m=1$では最初の和が0となり,初項にも合う.

    \noindent \textbf{【別解】}\\
    \,$\lfloor\sqrt k\rfloor$は$j^2\le k$を満たす正整数$j$の個数でもある.
    組$(j,k)$の条件$1\le k\le N$, $1\le j$, $j^2\le k$を,$j$ごとに数え直す.
    $k=j^2,\ldots,N$の$N-j^2+1$個から,
    \begin{align*}
      a_n&=\sum_{j=1}^{m}(N-j^2+1)\\
      &=m(N+1)-\frac{m(m+1)(2m+1)}6.
    \end{align*}
    群の和と,整数の組の数え直しが同じ答えを与えた.
\end{enumerate}

\subsubsection{余りと周期}\label{節：余りと周期}
整数$X$と正整数$m$に対して,
\[
  X-m\left\lfloor\frac Xm\right\rfloor
\]
は$X$を$m$で割った余りであり,$0,1,\ldots,m-1$のいずれかである.
\bookterm{floor-function}{床関数}の定義から,\textbf{値は0以上$m$未満の整数に限られる}.

$x\equiv y\pmod m$は,$x-y$が$m$の倍数であることを表す合同式である.
更新が$n$にも依存するなら,$a_n$と$n$の余りを組にして追う.
同じ組から次が一意に決まれば,戻った後は繰り返す.初項から周期的とは限らない.

\noindent \bookblocklink{example-gauss-role-3}{problem}{answer}{【例題】}\par\nobreak
\begin{enumerate}[label=(\arabic*),parsep=3pt,beginpenalty=10000,format=\bookitemlink{example-gauss-role-3}{problem}{answer}]
  \item \bookterm{sequence}{数列}を次のように定める.\bookterm{general}{一般項}と最小の正の周期を求めよ.
    \[
      a_1=1,\qquad
      a_{n+1}=X_n-4\left\lfloor\frac{X_n}{4}\right\rfloor,
    \]
    \[
      X_n=a_n^2+n^2+n\quad(n\ge1).
    \]
\end{enumerate}

\noindent \bookblocklink{example-gauss-role-3}{hint}{problem}{【方針】}\par\nobreak
\begin{enumerate}[label=(\arabic*),parsep=3pt,beginpenalty=10000,format=\bookitemlink{example-gauss-role-3}{hint}{problem}]
  \item 次の項は$X_n$を4で割った余りです.
    $(a_n,n\bmod4)$の組を追うと,同じ状態への戻りを確かめられます.
\end{enumerate}

\noindent \bookblocklink{example-gauss-role-3}{answer}{problem}{【解答】}\par\nobreak
\begin{enumerate}[label=(\arabic*),parsep=3pt,beginpenalty=10000,format=\bookitemlink{example-gauss-role-3}{answer}{problem}]
  \item \[
      a_n=
      \begin{cases}
        1 & (n\equiv0,1\pmod4),\\
        3 & (n\equiv2,3\pmod4).
      \end{cases}
    \]
    最小の正の周期は4である.
\end{enumerate}

\noindent \bookblocklink{example-gauss-role-3}{solution}{problem}{【解法】}\par\nobreak
\begin{enumerate}[label=(\arabic*),parsep=3pt,beginpenalty=10000,format=\bookitemlink{example-gauss-role-3}{solution}{problem}]
  \item 初項と更新から各項は整数で,第2項以降は$0\le a_n<4$である.
    \[
      a_{n+1}\equiv a_n^2+n(n+1)\pmod4
    \]
    となる.順に追うと,
    \[
      a_1,a_2,a_3,a_4,a_5=1,3,3,1,1.
    \]
    第1項と第5項では$(a_n,n\bmod4)=(1,1)$が一致する.
    同じ組から同じ次の組が決まるので,帰納的に$a_{n+4}=a_n$となり,【解答】を得る.
    第4項は値だけが同じで,添字の余りが違う.

    $a_1\ne a_2,a_3$と$a_2\ne a_5$より,1,2,3は周期でない.従って最小周期は4である.
\end{enumerate}

\subsubsection{境界と繰り上がり}\label{節：境界と繰り上がり}
正整数$n$と整数$m\ge2$に対して,
\[
  \left\lfloor\frac nm\right\rfloor-
  \left\lfloor\frac{n-1}{m}\right\rfloor
  =\begin{cases}
    1 & (n\text{が}m\text{の倍数}),\\
    0 & (\text{それ以外})
  \end{cases}
\]
となり,商が増える境界を検出する.
\,$\lfloor\sqrt n\rfloor-\lfloor\sqrt{n-1}\rfloor$なら平方数のときだけ1になる.
\textbf{条件を式で表し,\bookterm{recurrence}{漸化式}に組み込める.}

和の繰り上がりも検出できる.小数部分を$\{x\}=x-\lfloor x\rfloor$と書けば,
\begin{align*}
  &\lfloor x+y\rfloor-\lfloor x\rfloor-\lfloor y\rfloor\\
  &\qquad=\lfloor\{x\}+\{y\}\rfloor\in\{0,1\}.
\end{align*}
は小数部分の和が1以上のときだけ1となる.
\bookterm{floor-function}{床関数}を和に分配したときのずれは,繰り上がりを表す.

\noindent \bookblocklink{example-gauss-role-4}{problem}{answer}{【例題】}\par\nobreak
\begin{enumerate}[label=(\arabic*),parsep=3pt,beginpenalty=10000,format=\bookitemlink{example-gauss-role-4}{problem}{answer}]
  \item \bookterm{sequence}{数列}を次のように定める.\bookterm{general}{一般項}を求め,\bookterm{difference}{階差}の意味を説明せよ.
    \[
      a_0=0,\qquad a_n=a_{n-1}+d_n\quad(n\ge1),
    \]
    \[
      \begin{aligned}
        d_n={}&\left\lfloor\frac n4\right\rfloor
        -\left\lfloor\frac{n-1}{4}\right\rfloor\\
        &-\left\lfloor\frac n6\right\rfloor
        +\left\lfloor\frac{n-1}{6}\right\rfloor.
      \end{aligned}
    \]
\end{enumerate}

\noindent \bookblocklink{example-gauss-role-4}{hint}{problem}{【方針】}\par\nobreak
\begin{enumerate}[label=(\arabic*),parsep=3pt,beginpenalty=10000,format=\bookitemlink{example-gauss-role-4}{hint}{problem}]
  \item 各差を倍数かどうかの判定として読みます.
    和を取ると,途中の\bookterm{floor-function}{床関数}が打ち消し合います.
\end{enumerate}

\noindent \bookblocklink{example-gauss-role-4}{answer}{problem}{【解答】}\par\nobreak
\begin{enumerate}[label=(\arabic*),parsep=3pt,beginpenalty=10000,format=\bookitemlink{example-gauss-role-4}{answer}{problem}]
  \item \[
      a_n=\left\lfloor\frac n4\right\rfloor
      -\left\lfloor\frac n6\right\rfloor.
    \]
    \,$d_n$は4の倍数だけなら1,6の倍数だけなら$-1$,両方またはいずれでもないなら0である.
\end{enumerate}

\noindent \bookblocklink{example-gauss-role-4}{solution}{problem}{【解法】}\par\nobreak
\begin{enumerate}[label=(\arabic*),parsep=3pt,beginpenalty=10000,format=\bookitemlink{example-gauss-role-4}{solution}{problem}]
  \item \bookterm{floor-function}{床関数}の差の性質から,\,$d_n$は【解答】に述べた条件を表す.
    また,\,$a_n=\sum_{k=1}^{n}d_k$であり,
    \begin{align*}
      \sum_{k=1}^{n}\left(
      \left\lfloor\frac k4\right\rfloor-
      \left\lfloor\frac{k-1}{4}\right\rfloor\right)
      &=\left\lfloor\frac n4\right\rfloor,\\
      \sum_{k=1}^{n}\left(
      \left\lfloor\frac k6\right\rfloor-
      \left\lfloor\frac{k-1}{6}\right\rfloor\right)
      &=\left\lfloor\frac n6\right\rfloor.
    \end{align*}
    両式の差が【解答】となり,$n=0$にも合う.
    条件ごとの計算を,差の和で一度に処理した.
\end{enumerate}

\subsubsection{丸め誤差を使う}\label{節：丸め誤差を使う}
\,$a_{n+1}=\lfloor\lambda a_n\rfloor$では,$\lambda$の整数部分だけを取り出せない.そこで,
\[
  \varepsilon_n=\lambda a_n-a_{n+1},
  \qquad 0\le\varepsilon_n<1
\]
と置く.\textbf{捨てた部分の範囲から,次の整数部分を確定する.}

$\lambda^2$を$\lambda$と整数で表せる係数が手掛かりになる.
ただし,誤差の範囲を示せて初めて隣接3項間\bookterm{recurrence}{漸化式}へ帰着できる.

\noindent \bookblocklink{example-gauss-role-5}{problem}{answer}{【例題】}\par\nobreak
\begin{enumerate}[label=(\arabic*),parsep=3pt,beginpenalty=10000,format=\bookitemlink{example-gauss-role-5}{problem}{answer}]
  \item \bookterm{sequence}{数列}を次のように定める.\bookterm{floor}{ガウス記号}を含まない\bookterm{recurrence}{漸化式}に帰着させ,\bookterm{general}{一般項}を求めよ.
    \[
      a_1=1,\qquad
      a_{n+1}=\lfloor(1+\sqrt2)a_n\rfloor
      \quad(n\ge1).
    \]
\end{enumerate}

\noindent \bookblocklink{example-gauss-role-5}{hint}{problem}{【方針】}\par\nobreak
\begin{enumerate}[label=(\arabic*),parsep=3pt,beginpenalty=10000,format=\bookitemlink{example-gauss-role-5}{hint}{problem}]
  \item \,$\lambda=1+\sqrt2$とおき,\,$\lambda^2=2\lambda+1$を使います.
    整数から正の1未満の誤差を引く形にして,\bookterm{floor-function}{床関数}の値を確定しましょう.
\end{enumerate}

\noindent \bookblocklink{example-gauss-role-5}{answer}{problem}{【解答】}\par\nobreak
\begin{enumerate}[label=(\arabic*),parsep=3pt,beginpenalty=10000,format=\bookitemlink{example-gauss-role-5}{answer}{problem}]
  \item \[
      a_1=1,\quad a_2=2,\qquad
      a_{n+2}=2a_{n+1}+a_n-1.
    \]
    \,$\alpha=1+\sqrt2,\ \beta=1-\sqrt2$とおくと,
    \[
      a_n=\frac12+\frac{\alpha^n+\beta^n}{4}.
    \]
\end{enumerate}

\noindent \bookblocklink{example-gauss-role-5}{solution}{problem}{【解法】}\par\nobreak
\begin{enumerate}[label=(\arabic*),parsep=3pt,beginpenalty=10000,format=\bookitemlink{example-gauss-role-5}{solution}{problem}]
  \item \,$\lambda=1+\sqrt2>2$とおく.正整数$a_n$から次の項も正整数となる.
    また,\,$\lambda a_n$は無理数なので,
    \[
      \varepsilon_n=\lambda a_n-a_{n+1}
    \]
    は$0<\varepsilon_n<1$を満たす.
    \,$\lambda^2=2\lambda+1$を使うと,
    \begin{align*}
      \lambda a_{n+1}
      &=\lambda(\lambda a_n-\varepsilon_n)\\
      &=2a_{n+1}+a_n-(\lambda-2)\varepsilon_n.
    \end{align*}
    $M_n=2a_{n+1}+a_n$は整数で,$0<(\lambda-2)\varepsilon_n<1$だから,
    \[
      M_n-1<\lambda a_{n+1}<M_n.
    \]
    より$a_{n+2}=M_n-1$.\bookterm{initial}{初期条件}は$a_2=\lfloor1+\sqrt2\rfloor=2$である.

    定数項を消すため$b_n=a_n-\frac12$とおくと,
    \[
      b_{n+2}=2b_{n+1}+b_n,
      \qquad b_1=\frac12,\quad b_2=\frac32.
    \]
    \sectionref{節：隣接3項型}で学んだ\bookterm{characteristic}{特性方程式}は$x^2-2x-1=0$で,
    その解が$\alpha,\beta$である.
    \,$\alpha+\beta=2,\ \alpha^2+\beta^2=6$より,
    \[
      b_n=\frac{\alpha^n+\beta^n}{4}
    \]
    は初期値と更新を満たし,一意性から求める解となる.
    $a_n=b_n+\frac12$に戻せば【解答】を得る.
\end{enumerate}

\subsubsection{止まる値を調べる}\label{節：止まる値を調べる}
次の値が今の値と等しい\bookterm{fixed}{不動点}に達すると,\bookterm{sequence}{数列}は止まる.
丸めを外した実数の\bookterm{recurrence}{漸化式}とは,止まる値が異なることがある.

整数$K$を引いて,
\[
  b_{n+1}=\left\lfloor\frac{b_n}{2}\right\rfloor
\]
に帰着できる場合がある.
正整数$m$に対して,
\[
  \left\lfloor\frac{\lfloor x\rfloor}{m}\right\rfloor
  =\left\lfloor\frac xm\right\rfloor
\]
が成り立つ.$\lfloor x\rfloor=mq+r\ (0\le r<m)$と書けば,$x=mq+r+t\ (0\le t<1)$.
$0\le r+t<m$より,両方の\bookterm{floor-function}{床関数}が$q$となるからである.

\textbf{負数にも注意する.}$-1<x<0$なら$\lfloor x\rfloor=-1$である.

\noindent \bookblocklink{example-gauss-role-6}{problem}{answer}{【例題】}\par\nobreak
\begin{enumerate}[label=(\arabic*),parsep=3pt,beginpenalty=10000,format=\bookitemlink{example-gauss-role-6}{problem}{answer}]
  \item 非負整数$M$に対し,\bookterm{sequence}{数列}を次のように定める.
    \[
      a_1=M,\qquad
      a_{n+1}=\left\lfloor\frac{a_n+6}{2}\right\rfloor
      \quad(n\ge1).
    \]
    \bookterm{general}{一般項}と,最終的に一定になる値を求めよ.
\end{enumerate}

\noindent \bookblocklink{example-gauss-role-6}{hint}{problem}{【方針】}\par\nobreak
\begin{enumerate}[label=(\arabic*),parsep=3pt,beginpenalty=10000,format=\bookitemlink{example-gauss-role-6}{hint}{problem}]
  \item 6を引くと,\bookterm{floor-function}{床関数}によって半分にする\bookterm{recurrence}{漸化式}になります.
    初項が6より小さい場合は,負の数の\bookterm{floor-function}{床関数}を使って考えましょう.
\end{enumerate}

\noindent \bookblocklink{example-gauss-role-6}{answer}{problem}{【解答】}\par\nobreak
\begin{enumerate}[label=(\arabic*),parsep=3pt,beginpenalty=10000,format=\bookitemlink{example-gauss-role-6}{answer}{problem}]
  \item \[
      a_n=6+\left\lfloor\frac{M-6}{2^{n-1}}\right\rfloor.
    \]
    有限回の操作の後で,$M<6$なら5,$M\ge6$なら6に固定される.
\end{enumerate}

\noindent \bookblocklink{example-gauss-role-6}{solution}{problem}{【解法】}\par\nobreak
\begin{enumerate}[label=(\arabic*),parsep=3pt,beginpenalty=10000,format=\bookitemlink{example-gauss-role-6}{solution}{problem}]
  \item 初項からすべての項が非負整数となる.
    $b_n=a_n-6$とおくと,整数6を\bookterm{floor-function}{床関数}の外に出せるので,
    \begin{align*}
      b_{n+1}
      &=\left\lfloor\frac{b_n+12}{2}\right\rfloor-6\\
      &=\left\lfloor\frac{b_n}{2}\right\rfloor.
    \end{align*}
    \bookterm{floor-function}{床関数}の合成の性質と帰納法により,
    \[
      b_n=\left\lfloor\frac{M-6}{2^{n-1}}\right\rfloor.
    \]
    十分大きな$n$で,中身は$M\ge6$なら$[0,1)$,$M<6$なら$(-1,0)$に入る.
    床はそれぞれ0,$-1$なので,6または5に有限回で到達する.

    実際,整数$x$が\bookterm{fixed}{不動点}となる条件は,
    \[
      x\le\frac{x+6}{2}<x+1
      \iff 4<x\le6
    \]
    より5と6が\bookterm{fixed}{不動点}となる.初項による分岐は\sectionref{節：初項による違い}でも見る.
\end{enumerate}

\subsubsection{添字を分ける}\label{節：添字を分ける}
\bookterm{floor-function}{床関数}で,\textbf{参照する項}を指定する場合もある.
正整数$n$に対して,
\[
  \left\lfloor\frac{n+1}{2}\right\rfloor
  +\left\lfloor\frac n2\right\rfloor=n
\]
は$n$の二分割である.偶数なら同じ添字,奇数なら隣り合う添字になる.

偶奇で分けて床を外し,\bookterm{difference}{階差}を取る.
半分の添字との関係が出れば,$1,\ 2\le n<4,\ 4\le n<8,\ldots$の群で整理できる.
ただし,元と同じ添字を参照する式だけでは初項が決まらないこともある.

\noindent \bookblocklink{example-gauss-role-7}{problem}{answer}{【例題】}\par\nobreak
\begin{enumerate}[label=(\arabic*),parsep=3pt,beginpenalty=10000,format=\bookitemlink{example-gauss-role-7}{problem}{answer}]
  \item \bookterm{sequence}{数列}を次のように定める.\bookterm{general}{一般項}を求めよ.
    \BookMathLayout{            \[
        a_1=2,\qquad
        a_n=a_{\left\lfloor\frac{n+1}{2}\right\rfloor}
        +a_{\left\lfloor\frac n2\right\rfloor}+n-1
        \quad(n\ge2).
      \]}{            \[
        \begin{aligned}
          a_1&=2,\\
          a_n&=a_{\left\lfloor\frac{n+1}{2}\right\rfloor}
          +a_{\left\lfloor\frac n2\right\rfloor}+n-1
          \\              &\quad(n\ge2).
        \end{aligned}
      \]}
    必要ならば,$2^m\le n<2^{m+1}$を満たす整数$m\ge0$を用いてよい.
\end{enumerate}

\noindent \bookblocklink{example-gauss-role-7}{hint}{problem}{【方針】}\par\nobreak
\begin{enumerate}[label=(\arabic*),parsep=3pt,beginpenalty=10000,format=\bookitemlink{example-gauss-role-7}{hint}{problem}]
  \item 偶奇で式を分け,$b_n=a_{n+1}-a_n$を考えます.
    $b_{2k}$と$b_{2k+1}$が$b_k+1$になることを示し,群ごとに足しましょう.
\end{enumerate}

\noindent \bookblocklink{example-gauss-role-7}{answer}{problem}{【解答】}\par\nobreak
\begin{enumerate}[label=(\arabic*),parsep=3pt,beginpenalty=10000,format=\bookitemlink{example-gauss-role-7}{answer}{problem}]
  \item $2^m\le n<2^{m+1}$を満たす整数$m\ge0$に対して,
    \[
      a_n=(m+3)n-2^{m+1}+1.
    \]
\end{enumerate}

\noindent \bookblocklink{example-gauss-role-7}{solution}{problem}{【解法】}\par\nobreak
\begin{enumerate}[label=(\arabic*),parsep=3pt,beginpenalty=10000,format=\bookitemlink{example-gauss-role-7}{solution}{problem}]
  \item $n\ge2$では二つの添字は1以上$n-1$以下なので,\bookterm{sequence}{数列}は順に定義できる.
    偶奇で分けると,$k\ge1$について,
    \[
      \begin{aligned}
        a_{2k}&=2a_k+2k-1,\\
        a_{2k+1}&=a_{k+1}+a_k+2k.
      \end{aligned}
    \]
    $b_n=a_{n+1}-a_n$とおくと,両式の差から$b_{2k}=b_k+1$.
    また,$a_{2k+2}=2a_{k+1}+2k+1$を使うと$b_{2k+1}=b_k+1$である.
    $a_2=5$より$b_1=3$なので,
    \[
      b_{2k}=b_{2k+1}=b_k+1,\qquad b_1=3.
    \]

    $m=0$では$b_1=3$.次の群の添字を半分にすると前の群へ入り,値が1増える.
    群番号$m$について帰納的に,$2^m\le n<2^{m+1}$で$b_n=m+3$となる.

    従って,完全な群と最後の群を分けて足すと,
    \begin{align*}
      a_n={}&2+\sum_{j=0}^{m-1}2^j(j+3)\\
      &+(n-2^m)(m+3).
    \end{align*}
    ここで,
    \[
      \sum_{j=0}^{m-1}2^j=2^m-1,
    \]
    \[
      \sum_{j=0}^{m-1}j2^j=(m-2)2^m+2
    \]
    を使う(後者は和$T$に対する$2T-T$から得る).$m=0$でも空の和として成立する.
    代入すると,
    \[
      a_n=(m+3)n-2^{m+1}+1
    \]
    を得る.添字の分割が,階差の群\bookterm{sequence}{数列}へつながった.
\end{enumerate}

\subsubsection{桁を取り出す}\label{節：桁を取り出す}
正整数$b\ge2$に対し,非負整数$n$を$b$で割って,
\[
  q=\left\lfloor\frac nb\right\rfloor,
  \qquad r=n-b\left\lfloor\frac nb\right\rfloor
\]
とおくと,$n=bq+r,\ 0\le r<b$である.
$b$進法は,数字0から$b-1$を使い,位を$b$倍ずつにする表し方である.
$q$は最後の桁を除いた整数,$r$はその桁の数字となる.十進法なら一の位に当たる.
\textbf{床で添字を小さくし,余りで最後の桁を処理する.}

\noindent \bookblocklink{example-gauss-role-8}{problem}{answer}{【例題】}\par\nobreak
\begin{enumerate}[label=(\arabic*),parsep=3pt,beginpenalty=10000,format=\bookitemlink{example-gauss-role-8}{problem}{answer}]
  \item \bookterm{sequence}{数列}を次のように定める.
    \[
      a_0=0,\qquad
      a_n=a_{\left\lfloor\frac n3\right\rfloor}
      +n-3\left\lfloor\frac n3\right\rfloor
      \quad(n\ge1).
    \]
    $a_n$が$n$の3進表示とどのような関係を持つか説明し,$a_{14}$を求めよ.
\end{enumerate}

\noindent \bookblocklink{example-gauss-role-8}{hint}{problem}{【方針】}\par\nobreak
\begin{enumerate}[label=(\arabic*),parsep=3pt,beginpenalty=10000,format=\bookitemlink{example-gauss-role-8}{hint}{problem}]
  \item $n=3q+r\ (0\le r<3)$と書くと,\bookterm{recurrence}{漸化式}は$a_n=a_q+r$になります.
    $q$の3進表示に最後の桁$r$を付けると$n$になることを使いましょう.
\end{enumerate}

\noindent \bookblocklink{example-gauss-role-8}{answer}{problem}{【解答】}\par\nobreak
\begin{enumerate}[label=(\arabic*),parsep=3pt,beginpenalty=10000,format=\bookitemlink{example-gauss-role-8}{answer}{problem}]
  \item $a_n$は$n$の3進表示の各桁の数字の和である.
    $14=(112)_3$より,
    \[
      a_{14}=1+1+2=4.
    \]
\end{enumerate}

\noindent \bookblocklink{example-gauss-role-8}{solution}{problem}{【解法】}\par\nobreak
\begin{enumerate}[label=(\arabic*),parsep=3pt,beginpenalty=10000,format=\bookitemlink{example-gauss-role-8}{solution}{problem}]
  \item $n\ge1$では$0\le q=\lfloor\frac n3\rfloor<n$なので,初期値$a_0$から各項を定義できる.
    $r=n-3q$とおけば$r$は0,1,2のいずれかで,
    \[
      a_n=a_q+r.
    \]
    0の桁の和を0とする.$n=3q+r$の表示は,$q$の最後に$r$を付けたものなので,桁の和も$r$増える.
    $q<n$と$a_0=0$から,強い帰納法により各項は桁の和に一致する.

    実際,$14=3\cdot4+2,\ 4=3\cdot1+1,\ 1=3\cdot0+1$なので,
    \[
      a_{14}=a_4+2=a_1+3=a_0+4=4.
    \]
    値の意味を読むことも,\bookterm{sequence}{数列}を理解する方法である.
\end{enumerate}

\subsubsection{小数部分と回転}\label{節：小数部分と回転}
\,$\{x\}=x-\lfloor x\rfloor$は$x$の小数部分を表し,\,$0\le\{x\}<1$である.
例えば,\,$\{-0.2\}=0.8$となる.
一定の量を加えて整数部分を捨てれば,$[0,1)$へ戻る.
両端をつないで円周と見れば,\textbf{境界を越えても同じ割合で進む回転}である.
越えた回数は\bookterm{floor-function}{床関数}が記録する.周期は,進む割合が有理数かどうかで決まる.

\noindent \bookblocklink{example-gauss-role-9}{problem}{answer}{【例題】}\par\nobreak
\begin{enumerate}[label=(\arabic*),parsep=3pt,beginpenalty=10000,format=\bookitemlink{example-gauss-role-9}{problem}{answer}]
  \item \,$\alpha=\sqrt2-1$とし,\bookterm{sequence}{数列}を次のように定める.
    \[
      a_0=0,\qquad
      a_{n+1}=a_n+\alpha-\lfloor a_n+\alpha\rfloor
      \quad(n\ge0).
    \]
    \bookterm{general}{一般項}を求め,異なる添字の項が同じ値になるか調べよ.
\end{enumerate}

\noindent \bookblocklink{example-gauss-role-9}{hint}{problem}{【方針】}\par\nobreak
\begin{enumerate}[label=(\arabic*),parsep=3pt,beginpenalty=10000,format=\bookitemlink{example-gauss-role-9}{hint}{problem}]
  \item 整数部分を捨てると小数部分だけが残ることから,\,$n\alpha$の小数部分を予想します.
    一致する二項があれば,添字の差と$\alpha$の積が整数になることを使いましょう.
\end{enumerate}

\noindent \bookblocklink{example-gauss-role-9}{answer}{problem}{【解答】}\par\nobreak
\begin{enumerate}[label=(\arabic*),parsep=3pt,beginpenalty=10000,format=\bookitemlink{example-gauss-role-9}{answer}{problem}]
  \item \[
      a_n=n\alpha-\lfloor n\alpha\rfloor\quad(n\ge0).
    \]
    異なる添字の項は同じ値にならず,周期的ではない.
\end{enumerate}

\noindent \bookblocklink{example-gauss-role-9}{solution}{problem}{【解法】}\par\nobreak
\begin{enumerate}[label=(\arabic*),parsep=3pt,beginpenalty=10000,format=\bookitemlink{example-gauss-role-9}{solution}{problem}]
  \item \,$n=0$では予想した式は$a_0=0$を与える.
    \,$a_n=n\alpha-\lfloor n\alpha\rfloor$とすると,整数との和の性質から,
    \begin{align*}
      \lfloor a_n+\alpha\rfloor
      &=\lfloor(n+1)\alpha\rfloor-\lfloor n\alpha\rfloor,\\
      a_{n+1}
      &=(n+1)\alpha-\lfloor(n+1)\alpha\rfloor.
    \end{align*}
    より帰納的に【解答】を得る.各項は$[0,1)$に入る.

    \,$a_{n+p}=a_n\ (p\ge1)$とすると,
    \[
      p\alpha=\lfloor(n+p)\alpha\rfloor-\lfloor n\alpha\rfloor
    \]
    が整数となり,$\alpha$が無理数であることに矛盾する.従って異なる項は一致しない.

    同じ\bookterm{recurrence}{漸化式}で$\alpha$が有理数$\frac pq\ (0<p<q)$であり,\,$p,q$が互いに素なら,
    最小周期は$q$となる.$q\alpha$は整数で,$T\alpha$が整数となる正整数$T$は$q$の倍数だからである.
\end{enumerate}

\subsubsection{整数の個数を数える}\label{節：整数の個数を数える}
$X\ge0$なら$\lfloor X\rfloor$は$1\le k\le X$を満たす整数の個数である.
$\lfloor \dfrac{n}{m}\rfloor$なら,$1,\ldots,n$の中の$m$の倍数を数える.
\textbf{\bookterm{floor-function}{床関数}の和を,個数として読み直せる.}

座標がともに整数の点を格子点という.
$c>0$で,$\sum_{k=1}^{n}\lfloor ck\rfloor$は,$x=k$ごとに$1\le y\le ck$の整数$y$を数えた和である.
群\bookterm{sequence}{数列}の別解と同じく,整数の組を数える見方になる.

\noindent \bookblocklink{example-gauss-role-10}{problem}{answer}{【例題】}\par\nobreak
\begin{enumerate}[label=(\arabic*),parsep=3pt,beginpenalty=10000,format=\bookitemlink{example-gauss-role-10}{problem}{answer}]
  \item \bookterm{sequence}{数列}を次のように定める.
    \[
      a_0=0,\qquad
      a_n=a_{n-1}+\left\lfloor\frac{5n}{2}\right\rfloor
      \quad(n\ge1).
    \]
    \bookterm{general}{一般項}を求め,$a_n$をある範囲の格子点の個数として説明せよ.
\end{enumerate}

\noindent \bookblocklink{example-gauss-role-10}{hint}{problem}{【方針】}\par\nobreak
\begin{enumerate}[label=(\arabic*),parsep=3pt,beginpenalty=10000,format=\bookitemlink{example-gauss-role-10}{hint}{problem}]
  \item $\lfloor\frac{5k}{2}\rfloor=2k+\lfloor\frac k2\rfloor$と分けます.
    後者の和を偶奇で求め,\bookterm{floor-function}{床関数}が数えている整数$y$の範囲も確認しましょう.
\end{enumerate}

\noindent \bookblocklink{example-gauss-role-10}{answer}{problem}{【解答】}\par\nobreak
\begin{enumerate}[label=(\arabic*),parsep=3pt,beginpenalty=10000,format=\bookitemlink{example-gauss-role-10}{answer}{problem}]
  \item \[
      a_n=n(n+1)+\left\lfloor\frac{n^2}{4}\right\rfloor.
    \]
    これは整数$x,y$で,
    \[
      1\le x\le n,\qquad 1\le y\le\frac{5x}{2}
    \]
    を満たす格子点の個数である.
\end{enumerate}

\noindent \bookblocklink{example-gauss-role-10}{solution}{problem}{【解法】}\par\nobreak
\begin{enumerate}[label=(\arabic*),parsep=3pt,beginpenalty=10000,format=\bookitemlink{example-gauss-role-10}{solution}{problem}]
  \item \bookterm{difference}{階差}を足し,整数$2k$を\bookterm{floor-function}{床関数}の外に出すと,
    \begin{align*}
      a_n&=\sum_{k=1}^{n}\left\lfloor\frac{5k}{2}\right\rfloor\\
      &=n(n+1)+\sum_{k=1}^{n}\left\lfloor\frac k2\right\rfloor.
    \end{align*}
    後ろの和は,$n=2m$なら,
    \[
      0+1+1+2+2+\cdots+(m-1)+(m-1)+m=m^2,
    \]
    $n=2m+1$なら$m^2+m$である.
    従っていずれの場合も,
    \[
      \sum_{k=1}^{n}\left\lfloor\frac k2\right\rfloor
      =\left\lfloor\frac{n^2}{4}\right\rfloor
    \]
    となり,【解答】を得る.$n=0$にも合う.

    $x=k$にある格子点は,$1\le y\le\dfrac{5k}{2}$の$\lfloor\dfrac{5k}{2}\rfloor$個.
    $k=1,\ldots,n$について足した$a_n$が,【解答】の範囲の点の個数となる.
\end{enumerate}


\medskip
\noindent \textbf{【着眼点のまとめ】}\\
記号の位置と役割から方法を選ぶ.
余りと桁,群\bookterm{sequence}{数列}と数え上げは,同じ量を異なる側から捉えている.

\subsection{隣接4項型}\label{節：隣接4項型}
更新に使う項を増やして,
\[
  a_{n+3}+pa_{n+2}+qa_{n+1}+ra_n=0
\]
のような\bookterm{recurrence}{漸化式}を,$n\ge1$で考える.
初めの三項を自由に選ぶ三階の\bookterm{homogeneous}{同次}\bookterm{linear}{線形}\bookterm{recurrence}{漸化式}なので,解全体は3\bookterm{dimension}{次元}である(\sectionref*{節：初期条件と解の数}).
\sectionref*{節：補助数列を作る}と同じく,$d_{n+1}=\gamma d_n$を目標に,
\bookterm{auxiliary}{補助数列}$d_n=a_{n+2}+\alpha a_{n+1}+\beta a_n$を作る.
その条件は
\BookMathLayout{\[
    a_{n+3}+\alpha a_{n+2}+\beta a_{n+1}=\gamma \left(a_{n+2}+\alpha a_{n+1}+\beta a_{n}\right)
  \]}{
  \[
    \begin{aligned}
      a_{n+3}+\alpha a_{n+2}+\beta a_{n+1}\\
      {}=\gamma\left(a_{n+2}+\alpha a_{n+1}+\beta a_n\right)
    \end{aligned}
  \]
}
である.両辺を左側に集め,元の\bookterm{recurrence}{漸化式}と係数を比べると,
\[
  \alpha-\gamma=p,\qquad
  \beta-\gamma\alpha=q,\qquad
  -\gamma\beta=r
\]
となる.最初の二式から
\[
  \alpha=p+\gamma,\qquad
  \beta=q+p\gamma+\gamma^2
\]
なので,最後の式へ代入すると
\[
  \gamma^3+p\gamma^2+q\gamma+r=0
\]
を得る.実数解を一つ選べば,$\alpha,\beta$も決まる.\footnote{三次方程式は少なくとも一つ実数解を持つが,計算しやすい値を見つけられるかは別の問題である.}
例題では因数分解や\bookterm{difference}{階差}を手掛かりに選ぶ.
$d_1=a_3+\alpha a_2+\beta a_1$から$d_n$を求めると,
\[
  a_{n+2}+\alpha a_{n+1}+\beta a_n=d_1\gamma^{\,n-1}
\]
という3項間\bookterm{recurrence}{漸化式}が得られる.\footnote{$\gamma=0$の場合は$d_1$を初項として,$n\ge2$では$d_n=0$と読む.
  また,\sectionref{節：定数の付加項がある場合}で扱ったのは右辺が定数$r$の場合であり,ここでは右辺が指数関数$d_1\gamma^{n-1}$になる点が異なる.}
$d_n$を求めた後,$a_n$へ戻す.
\[
  \begin{array}{c|c}
    \text{求める量}&\text{使う初期値}\\ \hline
    d_n&d_1\\
    a_n&a_1,a_2
  \end{array}
\]
$a_3$の情報は$d_1$に入っている.
得られた式は\sectionref{節：階差等比複合型}の考え方で,右辺と同じ更新をする量を引いて解ける場合がある.
例えば第1問では,$p=-3,q=3,r=-1$なので上の方程式は$(\gamma-1)^3=0$となり,
$\gamma=1,\alpha=-2,\beta=1$を得る.
こうして$d_n=a_{n+2}-2a_{n+1}+a_n$が決まる.
元の式を\mbox{三\bookterm{difference}{階差}}と見ても,同じ量が得られる.

\noindent \bookblocklink{example-four-term-1}{problem}{answer}{【例題】}\par\nobreak
\begin{enumerate}[label=(\arabic*),parsep=3pt,format=\bookitemlink{example-four-term-1}{problem}{answer}]
  \item $\displaystyle a_1=1,\,a_2=2,\,a_3=4$\\
    $a_{n+3}=3a_{n+2}-3a_{n+1}+a_n\quad(n\ge1)$
  \item $\displaystyle a_1=2,\,a_2=1,\,a_3=5$\\
    $a_{n+3}=2a_{n+2}+a_{n+1}-2a_n\quad(n\ge1)$
  \item $\displaystyle a_1=1,\,a_2=4,\,a_3=11$\\
    $a_{n+3}=5a_{n+2}-7a_{n+1}+3a_n\quad(n\ge1)$
\end{enumerate}

\noindent\bookblocklink{example-four-term-1}{hint}{problem}{【方針】}\par\nobreak
\begin{enumerate}[label=(\arabic*),parsep=3pt,format=\bookitemlink{example-four-term-1}{hint}{problem}]
  \item $d_n=a_{n+2}-2a_{n+1}+a_n$とおき,\bookterm{difference}{階差}をさらにとると一定になることを利用します.
  \item $d_n=a_{n+2}-a_{n+1}-2a_n$とおくと$d_{n+1}=d_n$となります.\bookterm{initial}{初期条件}から$d_n=0$を導き,3項間\bookterm{recurrence}{漸化式}として解きます.
  \item $d_n=a_{n+2}-2a_{n+1}+a_n$とおくと$d_{n+1}=3d_n$となります.得られた右辺を満たす\bookterm{sequence}{数列}を引き,\bookterm{difference}{階差}が一定の形に帰着させます.
\end{enumerate}

\noindent\bookblocklink{example-four-term-1}{answer}{problem}{【解答】}\par\nobreak
\begin{enumerate}[label=(\arabic*),parsep=3pt,format=\bookitemlink{example-four-term-1}{answer}{problem}]
  \item $\displaystyle a_n=\frac{n^2-n+2}{2}$
  \item $\displaystyle a_n=2^{n-1}+(-1)^{n-1}$
  \item $\displaystyle a_n=n-1+3^{n-1}$
\end{enumerate}

\noindent\bookblocklink{example-four-term-1}{solution}{problem}{【解法】}\par\nobreak
\begin{enumerate}[label=(\arabic*),parsep=3pt,format=\bookitemlink{example-four-term-1}{solution}{problem}]
  \item $d_n=a_{n+2}-2a_{n+1}+a_n$と置くと,
    \[
      d_{n+1}-d_n=a_{n+3}-3a_{n+2}+3a_{n+1}-a_n=0
    \]
    より$d_n=d_1=1$である.
    $e_n=a_{n+1}-a_n$と置けば,$e_{n+1}-e_n=1,e_1=1$より$e_n=n$.従って,
    \[
      a_n=a_1+\sum_{k=1}^{n-1}e_k
      =1+\sum_{k=1}^{n-1}k
      =\frac{n^2-n+2}{2}
    \]
    を得る.
  \item
    \[
      d_n=a_{n+2}-a_{n+1}-2a_n
    \]
    とおくと,
    \[
      d_{n+1}-d_n=a_{n+3}-2a_{n+2}-a_{n+1}+2a_n=0
    \]
    である.また,$d_1=5-1-2\cdot2=0$だから$d_n=0$であり,
    \[
      a_{n+2}-a_{n+1}-2a_n=0
    \]
    となる.\bookterm{characteristic}{特性方程式}$x^2-x-2=0$の解は$2,-1$なので,
    \[
      a_n=A\,2^{n-1}+B(-1)^{n-1}
    \]
    とおける.\bookterm{initial}{初期条件}より$A+B=2,\,2A-B=1$であるから$A=B=1$を得る.
  \item
    \[
      d_n=a_{n+2}-2a_{n+1}+a_n
    \]
    とおくと,与えられた\bookterm{recurrence}{漸化式}から$d_{n+1}=3d_n$である.\bookterm{initial}{初期条件}より$d_1=11-2\cdot4+1=4$なので,
    \[
      a_{n+2}-2a_{n+1}+a_n=4\cdot3^{n-1}
    \]
    となる.一方,$3^{n-1}$の2\bookterm{difference}{階差}も$4\cdot3^{n-1}$である.そこで$b_n=a_n-3^{n-1}$とおくと,
    \[
      b_{n+2}-2b_{n+1}+b_n=0
    \]
    となる.$b_1=0,\,b_2=1$より$b_n=n-1$であるから,
    \[
      a_n=b_n+3^{n-1}=n-1+3^{n-1}
    \]
    を得る.
\end{enumerate}

\subsection{特性根が虚数の場合}\label{節：特性根が虚数の場合}

隣接3\,項間\bookterm{recurrence}{漸化式}(\sectionref{節：隣接3項型})の式
$a_n=A\alpha^{n-1}+B\beta^{n-1}$は,異なる\bookterm{root}{特性根}が虚数でも成り立つ.
ここでは\bookterm{trigonometric}{三角関数}を使い,実数の式へ直す.

\[
  a_{n+2}+pa_{n+1}+qa_n=0\quad(n\ge1)
\]
の$p,q,a_1,a_2$を実数とする.\bookterm{characteristic}{特性方程式}$x^2+px+q=0$が
$p^2-4q<0$を満たすと,根は互いに\bookterm{conjugate}{共役}な$\alpha,\overline\alpha$となる.
例えば$a_{n+2}=2a_{n+1}-2a_n$では根が$1\pm i$なので,
\[
  a_n=A(1+i)^{n-1}+B(1-i)^{n-1}
\]
と書ける.この式の実数性を,\bookterm{polar}{極形式}から読み直そう.
初めの二項も更新の係数も実数だから,実際の各項も実数である.
複素数を使った二つの項の中で,虚部が打ち消し合う形を取り出したい.

虚部が正の根$\alpha$の距離を$\rho$,角度を$\theta$とする.
この角度が\bookterm{argument}{偏角}であり,$\rho>0$, $0<\theta<\pi$である.\footnote{$e^{it}=\cos t+i\sin t$という表記を使う.}
\[
  \alpha=\rho e^{i\theta},\qquad
  \overline\alpha=\rho e^{-i\theta}
\]
特性根を一回掛けると,距離は$\rho$倍,向きは$\theta$だけ進む.
従って,$t=(n-1)\theta$と置けば
\[
  \begin{aligned}
    \alpha^{n-1}&=\rho^{n-1}(\cos t+i\sin t),\\
    \overline\alpha^{\,n-1}&=\rho^{n-1}(\cos t-i\sin t)
  \end{aligned}
\]
となる.
\BookFigComplexRoots

根の二つの冪は共役で,余弦の部分は同じ,正弦につく虚数の符号は反対になっている.
$A,B$も\bookterm{conjugate}{共役}にして虚部を消す.
和$A+B$と$i(A-B)$が実数になるよう,その値を$C,D$と書く.
\par\noindent\begin{minipage}{\linewidth}
  実数$C,D$を用いて,
  \[
    A=\frac{C-iD}{2},\qquad B=\frac{C+iD}{2}
  \]
  と置く.
\end{minipage}\par
\BookMathLayout{\[
    \begin{aligned}
      &A\alpha^{n-1}+B\overline\alpha^{\,n-1}\\
      &\quad=\rho^{n-1}\{(A+B)\cos t\\
      &\qquad+i(A-B)\sin t\}\\
      &\quad=\rho^{n-1}(C\cos t+D\sin t)
    \end{aligned}
  \]}{\[
    \begin{aligned}
      &A\alpha^{n-1}+B\overline\alpha^{\,n-1}\\
      &\quad=\rho^{n-1}\{(A+B)\cos t\\
      &\qquad+i(A-B)\sin t\}\\
      &\quad=\rho^{n-1}(C\cos t+D\sin t)
    \end{aligned}
  \]}
より,
\BookMathLayout{\[
    a_n=\rho^{n-1}\{C\cos((n-1)\theta)+D\sin((n-1)\theta)\}
  \]}{\[
    \begin{aligned}
      a_n=\rho^{n-1}\{&C\cos((n-1)\theta)\\
      &+D\sin((n-1)\theta)\}
    \end{aligned}
  \]}
を得る.逆に実数列なら$B=\overline A$が必要なことは,
\hyperref[解釈：虚数解と基本解]{「二つの基本解を選び直す」}で確かめる.

\bookstep{初めの二項を合わせる}
$n=1,2$を代入すると,
\[
  a_1=C,\qquad a_2=\rho(C\cos\theta+D\sin\theta)
\]
なので,
\BookMathLayout{\[
    C=a_1,\qquad D=\cfrac{\frac{a_2}{\rho}-a_1\cos\theta}{\sin\theta}
  \]}{\[
    \begin{aligned}
      C&=a_1,\\
      D&=\cfrac{\frac{a_2}{\rho}-a_1\cos\theta}{\sin\theta}
    \end{aligned}
  \]}
と決まる.$\sin\theta\ne0$だから割れる.
根から$\rho,\theta$を読み,\bookterm{initial}{初期条件}を合わせればよい.

\noindent \bookblocklink{example-complex-1}{problem}{answer}{【例題】}\par\nobreak
\begin{enumerate}[label=(\arabic*),parsep=3pt,format=\bookitemlink{example-complex-1}{problem}{answer}]
  \item $\displaystyle a_1=1,\,a_2=0,\quad a_{n+2}=-a_n\quad(n\ge1)$
  \item \BookMathLayout{$\displaystyle a_1=0,\,a_2=1,\quad a_{n+2}=2a_{n+1}-2a_n\quad(n\ge1)$}{$\displaystyle a_1=0,\,a_2=1$\\ $\displaystyle a_{n+2}=2a_{n+1}-2a_n$\\ $\displaystyle(n\ge1)$}
  \item $\displaystyle a_1=1,\,a_2=2,\quad a_{n+2}=-4a_n\quad(n\ge1)$
  \item \BookMathLayout{$\displaystyle a_1=1,\,a_2=0,\quad a_{n+2}=2a_{n+1}-4a_n\quad(n\ge1)$}{$\displaystyle a_1=1,\,a_2=0$\\ $\displaystyle a_{n+2}=2a_{n+1}-4a_n$\\ $\displaystyle(n\ge1)$}
\end{enumerate}

\noindent\bookblocklink{example-complex-1}{hint}{problem}{【方針】}\par\nobreak
\begin{enumerate}[label=(\arabic*),parsep=3pt,format=\bookitemlink{example-complex-1}{hint}{problem}]
  \item 初項から数項を計算して\bookterm{periodic}{周期性}を見つけ,\bookterm{general}{一般項}を予想します.その後,\bookterm{induction}{数学的帰納法}で証明します.
  \item \bookterm{root}{特性根}を\bookterm{polar}{極形式}で表し,$\rho=\sqrt2,\,\theta=\dfrac{\pi}{4}$を用いて\bookterm{general}{一般項}を置きます.
  \item \bookterm{root}{特性根}は$\pm2i$です.$\rho=2,\,\theta=\dfrac{\pi}{2}$として\bookterm{initial}{初期条件}を代入します.
  \item \bookterm{root}{特性根}を\bookterm{polar}{極形式}で表し,$\rho=2,\,\theta=\dfrac{\pi}{3}$として係数を決定します.
\end{enumerate}

\noindent\bookblocklink{example-complex-1}{answer}{problem}{【解答】}\par\nobreak
\begin{enumerate}[label=(\arabic*),parsep=3pt,format=\bookitemlink{example-complex-1}{answer}{problem}]
  \item $\displaystyle a_n=\cos\frac{(n-1)\pi}{2}$
  \item $\displaystyle a_n=(\sqrt2)^{n-1}\sin\frac{(n-1)\pi}{4}$
  \item $\displaystyle a_n=2^{n-1}\left(\cos\frac{(n-1)\pi}{2}+\sin\frac{(n-1)\pi}{2}\right)$
  \item \BookMathLayout{$\displaystyle a_n=2^{n-1}\left(\cos\frac{(n-1)\pi}{3}-\frac1{\sqrt3}\sin\frac{(n-1)\pi}{3}\right)$}{$\displaystyle a_n=2^{n-1}\Bigl(\cos\frac{(n-1)\pi}{3}$\\
      $\displaystyle\hspace{3zw}-\frac1{\sqrt3}\sin\frac{(n-1)\pi}{3}\Bigr)$}
\end{enumerate}

\noindent\bookblocklink{example-complex-1}{solution}{problem}{【解法】}\par\nobreak
\begin{enumerate}[label=(\arabic*),parsep=3pt,format=\bookitemlink{example-complex-1}{solution}{problem}]
  \item 初めの並びは$1,0,-1,0,1,0,-1,0,\ldots$なので,
    $a_n=\cos\frac{(n-1)\pi}{2}$と予想する.
    $n=1,2$で初期値に合い,$a_k=\cos\frac{(k-1)\pi}{2}$なら,
    \[
      \begin{aligned}
        a_{k+2}&=-a_k=\cos\left(\frac{(k-1)\pi}{2}+\pi\right)\\
        &=\cos\frac{(k+1)\pi}{2}.
      \end{aligned}
    \]
    従って,奇数番目と偶数番目それぞれの\bookterm{induction}{数学的帰納法}で全項について成り立つ.
  \item \bookterm{characteristic}{特性方程式}$x^2-2x+2=0$の根は$1\pm i$.
    $\rho=\sqrt2,\theta=\dfrac{\pi}{4}$である.
    $a_1=0$から$C=0$,$a_2=1$から$D=1$と決まり,【解答】を得る.
  \item $x^2+4=0$の根は$\pm2i$なので,$\rho=2,\theta=\dfrac{\pi}{2}$.
    $a_1=1$から$C=1$,$a_2=2$から$2D=2$,従って$D=1$となる.
  \item $x^2-2x+4=0$の根は$1\pm\sqrt3\,i$なので,$\rho=2,\theta=\dfrac{\pi}{3}$.
    $a_1=1$から$C=1$,$a_2=0$から
    \[
      0=2\left(\frac12+D\frac{\sqrt3}{2}\right)
    \]
    となり,$D=-\dfrac{1}{\sqrt3}$を得る.
\end{enumerate}

\bookstep{加法定理による別解}\label{解釈：虚数解と加法定理}
複素数の冪を直す代わりに,更新そのものから三角関数を見つけよう.
解と係数の関係より,
\[
  p=-(\alpha+\overline\alpha)=-2\rho\cos\theta,
  \qquad q=\alpha\overline\alpha=\rho^2.
\]
そこで$b_n=\dfrac{a_n}{\rho^{n-1}}$と置き,元の式を$\rho^{n+1}$で割ると,
\[
  b_{n+2}+b_n=2\cos\theta\,b_{n+1}
\]
を得る.これは\bookterm{addition}{加法定理}
\[
  \begin{aligned}
    \cos(x+\theta)+\cos(x-\theta)&=2\cos\theta\cos x,\\
    \sin(x+\theta)+\sin(x-\theta)&=2\cos\theta\sin x
  \end{aligned}
\]
と同じ形である.$x=n\theta$を代入すれば,
$\cos((n-1)\theta)$も$\sin((n-1)\theta)$もこの更新を満たす.
従ってその組み合わせ
\[
  h_n=C\cos((n-1)\theta)+D\sin((n-1)\theta)
\]
も解になる.初めの二項は$b_1=a_1,b_2=\dfrac{a_2}{\rho}$なので,
$C=b_1$, $D=\dfrac{b_2-b_1\cos\theta}{\sin\theta}$と合わせる.

$b_1=h_1,b_2=h_2$であり,$b_n=h_n,b_{n+1}=h_{n+1}$なら,
\[
  b_{n+2}=2\cos\theta\,h_{n+1}-h_n=h_{n+2}.
\]
\textbf{加法定理で候補の更新を確かめ,帰納法で全項の一致を得た.}
最後に$a_n=\rho^{n-1}b_n$へ戻せば,前と同じ一般項になる.

\bookstep{振幅・位相・周期}
$C=D=0$なら零数列である.それ以外では$R=\sqrt{C^2+D^2}>0$とし,
\[
  \cos\varphi=\frac CR,\qquad \sin\varphi=\frac DR
\]
を満たす角$\varphi$を選ぶ.加法定理から,
\[
  a_n=R\rho^{n-1}\cos((n-1)\theta-\varphi)
\]
とまとまる.$\varphi$は波の位置をずらす角度(位相),$\theta$は一回で進む角度である.
振幅$R\rho^{n-1}$は,
\[
  \begin{cases}
    0<\rho<1 & \text{減衰する},\\
    \rho=1 & \text{一定である},\\
    \rho>1 & \text{増大する}.
  \end{cases}
\]
これは$|a_n|\le R\rho^{n-1}$という上限であり,各項の絶対値が同じ割合で変わるという意味ではない.

三角関数の式でも,整数添字で周期的とは限らない.
$\rho=1$かつ$\dfrac{\theta}{2\pi}$が有理数なら,$T\theta=2k\pi$となる正整数$T$と整数$k$が存在する.
角度が一周の整数倍だけ進むので,$a_{n+T}=a_n$となる.
零数列を除けば,この二条件は必要でもある.次に独立性から確かめる.

\bookstep{二つの基本解を選び直す}\phantomsection\label{解釈：虚数解と基本解}
\sectionref*{節：補助数列を作る}では,目標の更新に合う量を作った.
$b_n=\dfrac{a_n}{\rho^{n-1}}$も,根の絶対値から\bookterm{normalize}{正規化}の倍率を選び,
係数の$\rho,\rho^2$を取り除く操作である.残る更新が加法定理になった.

\sectionref*{節：初期条件と解の数}の見方では,解を表す\bookterm{basis}{基底}を選び直している.
$p^2-4q<0$から$q>0$で,根$\alpha,\overline\alpha$は異なる非零の数である.
二つの冪$\alpha^{n-1},\overline\alpha^{\,n-1}$の\bookterm{combination}{線形結合}が零数列なら,初めの二項から
\[
  U+V=0,\qquad U\alpha+V\overline\alpha=0
\]
となる.$\alpha\ne\overline\alpha$より$U=V=0$なので,\bookterm{independent}{一次独立}である.
同じ連立式で任意の初期二項を合わせられ,更新の一意性から全項も一致する.
従って二つの冪は複素数係数で考えた解全体の基底であり,解全体は2\bookterm{dimension}{次元}となる.

実数解なら$a_n=\overline{a_n}$より,
\[
  A\alpha^{n-1}+B\overline\alpha^{\,n-1}
  =\overline B\alpha^{n-1}+\overline A\,\overline\alpha^{\,n-1}.
\]
係数の一意性から$A=\overline B,B=\overline A$が必要である.
一方,実数解$\rho^{n-1}\cos((n-1)\theta)$と$\rho^{n-1}\sin((n-1)\theta)$の初めの二項は,
それぞれ$(1,\rho\cos\theta)$と$(0,\rho\sin\theta)$である.
$\rho\sin\theta\ne0$なので,こちらは実数係数で考えた解全体の基底となる.
虚部を消す計算は,実数の解を表す基底への変更だった.

\bookterm{periodic}{周期性}の必要条件も示せる.周期$T$があるなら,
\[
  A(\alpha^T-1)\alpha^{n-1}
  +\overline A(\overline\alpha^{\,T}-1)\overline\alpha^{\,n-1}=0.
\]
独立性から両係数が零となる.非零の実数解では$A,\overline A\ne0$なので,
$\alpha^T=\overline\alpha^{\,T}=1$.従って$\rho=1$,$\dfrac{\theta}{2\pi}$は有理数である.

\bookstep{補足：指数と初項の番号}
主な導出では$a_1,a_2$と指数$n-1$を使った.指数$n$の表示と,第0項から始める表示も区別しておこう.

\bookstep{同じ数列を指数$n$で表す}
非零の根なら$A\alpha^{n-1}=\left(\dfrac{A}{\alpha}\right)\alpha^n$と係数を取り直せる.
$\widetilde A=\dfrac{A}{\alpha}$, $\widetilde B=\dfrac{B}{\overline\alpha}$とすれば,
\[
  \begin{aligned}
    a_n&=\widetilde A\alpha^n+\widetilde B\overline\alpha^{\,n}\\
    &=\rho^n(c\cos n\theta+d\sin n\theta).
  \end{aligned}
\]
$c,d$は$C,D$と別の係数である.
$a_1=\rho(c\cos\theta+d\sin\theta)$,
$a_2=\rho^2(c\cos2\theta+d\sin2\theta)$から決まり,対応は
\BookMathLayout{\[
    c=\frac{C\cos\theta-D\sin\theta}{\rho},\qquad
    d=\frac{C\sin\theta+D\cos\theta}{\rho}
  \]}{\[
    \begin{aligned}
      c&=\frac{C\cos\theta-D\sin\theta}{\rho},\\
      d&=\frac{C\sin\theta+D\cos\theta}{\rho}
    \end{aligned}
  \]}
となる.変えたのは表示と係数であり,項の番号ではない.

\bookstep{第0項から始める場合}
初期条件を$a_0,a_1$とし,$n\ge0$で更新する場合は,別の係数$C_0,D_0$を使う.
指数$n$の形に$n=0,1$を代入して,
\[
  C_0=a_0,\qquad
  D_0=\cfrac{\frac{a_1}{\rho}-a_0\cos\theta}{\sin\theta}
\]
と決まる.例えば$a_{n+2}=2a_{n+1}-2a_n$,$a_0=1,a_1=2$では,
根$1\pm i$から$\rho=\sqrt2,\theta=\dfrac{\pi}{4}$.
$C_0=1$, $2=\dfrac{\sqrt2(C_0+D_0)}{\sqrt2}=1+D_0$より$D_0=1$なので,
\[
  a_n=(\sqrt2)^n\left(\cos\frac{n\pi}{4}+\sin\frac{n\pi}{4}\right)
\]
となる.

\subsection{非隣接2項型}\label{節：非隣接2項型}

非隣接2項間\bookterm{recurrence}{漸化式}
\[
  a_{n+2}=f(a_n)\qquad(n\ge1)
\]
では,$a_1$から奇数番目,$a_2$から偶数番目が決まる.
$u_k=a_{2k-1},v_k=a_{2k}$と部分列に分ける.
\[
  \begin{array}{c|cccc}
    k&1&2&3&\cdots\\ \hline
    u_k&a_1&a_3&a_5&\cdots\\
    v_k&a_2&a_4&a_6&\cdots
  \end{array}
\]
$n$が2進むと$k$は1進み,$u_{k+1}=f(u_k),v_{k+1}=f(v_k)$となる.
最後に奇数$n$では$k=\dfrac{n+1}{2}$,偶数$n$では$k=\dfrac{n}{2}$で戻す.

また,$a_{n+2}=qa_n$では,$a_{n+2}=0\cdot a_{n+1}+qa_n$と見て\bookterm{characteristic}{特性方程式}を使うこともできる.

\noindent \bookblocklink{example-nonadjacent-1}{problem}{answer}{【例題】}\par\nobreak
\begin{enumerate}[label=(\arabic*),parsep=3pt,format=\bookitemlink{example-nonadjacent-1}{problem}{answer}]
  \item $\displaystyle a_1=2,\,a_2=3,\,a_{n+2}=a_n^2$
  \item $\displaystyle a_1=1,\,a_2=2,\,a_{n+2}=2a_n+1$
  \item $\displaystyle a_1=1,\,a_2=3,\,a_{n+2}=4a_n$
\end{enumerate}

\noindent\bookblocklink{example-nonadjacent-1}{hint}{problem}{【方針】}\par\nobreak
\begin{enumerate}[label=(\arabic*),parsep=3pt,format=\bookitemlink{example-nonadjacent-1}{hint}{problem}]
  \item 奇数番目と偶数番目を別々の\bookterm{sequence}{数列}とみなし,それぞれの二乗関係を調べます.
  \item $b_n=a_n+1$とおいて定数項を消してから,奇数番目と偶数番目に分けます.
  \item $a_{n+2}=0\cdot a_{n+1}+4a_n$と見て,\bookterm{characteristic}{特性方程式}を解きます.
\end{enumerate}

\noindent\bookblocklink{example-nonadjacent-1}{answer}{problem}{【解答】}\par\nobreak
\begin{enumerate}[label=(\arabic*),parsep=3pt,format=\bookitemlink{example-nonadjacent-1}{answer}{problem}]
  \item $\displaystyle a_{2k-1}=2^{\,2^{k-1}},\quad a_{2k}=3^{\,2^{k-1}}\quad(k\ge1)$
  \item $\displaystyle a_{2k-1}=2^k-1,\quad a_{2k}=3\cdot2^{k-1}-1\quad(k\ge1)$
  \item $\displaystyle a_n=\frac54\,2^{n-1}-\frac14\,(-2)^{n-1}$
\end{enumerate}

\noindent\bookblocklink{example-nonadjacent-1}{solution}{problem}{【解法】}\par\nobreak
\begin{enumerate}[label=(\arabic*),parsep=3pt,format=\bookitemlink{example-nonadjacent-1}{solution}{problem}]
  \item $u_k=a_{2k-1},\,v_k=a_{2k}$とおくと,
    \[
      u_1=2,\quad v_1=3,\qquad
      u_{k+1}=u_k^2,\quad v_{k+1}=v_k^2
    \]
    となる.したがって,それぞれ順に2乗していくので,
    \[
      u_k=2^{\,2^{k-1}},\qquad v_k=3^{\,2^{k-1}}
    \]
    となり,$u_k=a_{2k-1},v_k=a_{2k}$で戻す.
  \item $b_n=a_n+1$とおくと,
    \[
      b_{n+2}=a_{n+2}+1=2a_n+2=2b_n
    \]
    であり,$b_1=2,\,b_2=3$である.
    よって,奇数番目と偶数番目で分けると,
    \[
      b_{2k-1}=2^k,\qquad b_{2k}=3\cdot2^{k-1}
    \]
    となる.$a_n=b_n-1$に戻せば【解答】を得る.
  \item \bookterm{characteristic}{特性方程式}は
    \[
      r^2-4=0
    \]
    であり,解は$r=2,-2$である.したがって,
    \[
      a_n=A\,2^{n-1}+B\,(-2)^{n-1}
    \]
    とおける.$a_1=1,\,a_2=3$を代入すると,
    \[
      A+B=1,\qquad 2A-2B=3
    \]
    より$A=\dfrac{5}{4},B=-\dfrac{1}{4}$と決まり,【解答】を得る.
\end{enumerate}

\subsubsection{一周期分をまとめる}\label{節：一周期分をまとめる}
更新の係数が奇数回と偶数回で変わる場合も,二回分をまとめられる.

\bookstep{例題：二回分をまとめる}
\[
  a_1=1,\qquad
  a_{n+1}=\begin{cases}
    2a_n+1&(n\text{が奇数}),\\
    3a_n+2&(n\text{が偶数})
  \end{cases}
\]
奇数番目を$x_k=a_{2k-1}$と置くと,$a_{2k}=2x_k+1$であり,
次の更新まで続けて
\[
  x_{k+1}=3(2x_k+1)+2=6x_k+5
\]
を得る.従って$x_1=1$から$x_k=2\cdot6^{k-1}-1$であり,
\BookMathLayout{\[
    a_{2k-1}=2\cdot6^{k-1}-1,\qquad
    a_{2k}=4\cdot6^{k-1}-1\quad(k\ge1)
  \]}{\[
    \begin{aligned}
      a_{2k-1}&=2\cdot6^{k-1}-1,\\
      a_{2k}&=4\cdot6^{k-1}-1\quad(k\ge1)
    \end{aligned}
  \]}
となる.偶数番目は更新の途中から戻した.
付加項も次の倍率を受けるため,二回分は$1+2$でなく$3\cdot1+2=5$である.

\bookstep{一周期への拡張}
$a_{n+1}=p_na_n+q_n$の係数が周期$L$をもつなら,
$x_k=a_{1+kL}$ $(k\ge0)$と置き,$x_0=a_1$から一周期分を合成する.
\[
  x_{k+1}=Px_k+Q
\]
の$P$は倍率の積,$Q$は各付加項にその後の倍率を掛けた寄与の和である.\footnote{最初の一周期の係数を$p_1,\ldots,p_L$, $q_1,\ldots,q_L$とすると,
  $P=\prod_{j=1}^L p_j$, $Q=\sum_{j=1}^L q_j\prod_{i=j+1}^L p_i$である.空積は$1$とする.
  この表示は倍率に$0$を含む場合も成り立つ.}
途中の項は元の一回ずつの更新から戻せる.
$P=0$なら$x_k=Q$ $(k\ge1)$,$P=1$なら$x_k=x_0+kQ$である.
それ以外は,\sectionref*{節：特性方程式型}と同じく不動点からの差を取る.

\begin{bookpoint}
  \textbf{更新の間隔を選ぶ}\par
  係数の一周期分をまとめる.係数の周期性は,項の値の周期性とは異なる.
\end{bookpoint}
付録の\sectionref*{節：二回分の更新}では,同じ計算を行列の積で読む.
符号も変わる例は\bookproblemref{practice-connections}{4}{接続演習問4}へ.

